%% ============================================================================
%% revised_introduction.tex
%%
%% Replacement for Section 1 of paper.tex (original lines 162--267).
%% Addresses referee:
%%   - resolves the loose use of "functor": adopts dual interpretation,
%%     defining the constellation map and pinning down a precise category
%%     (Def 1.1) on which it is a functor; the categorical claim is
%%     factual and minimal.
%%   - new chi_w-vs.-edge-indicial paragraph (sec:intro-chi-w).
%%   - new scope-of-analytic-lifting paragraph (sec:intro-analytic-scope).
%%   - explicit new-vs.-classical statement in each major subsection.
%%   - cleaned roadmap; no dataset-internal identifiers in the prose.
%% ============================================================================

\section{Introduction}\label{sec:intro}

\subsection{The constellation map}\label{ssec:intro-functor}

To a scalar linear operator $L$ with a single irregular singularity at
$x = \infty$, the WKB ansatz on the dominant Newton edge attaches a
finite multiset of complex numbers, the \emph{exponent constellation}
\[
   \Cl(L) \;\subset\; \mathbb{C}.
\]
The present paper is about the combinatorial geometry of the
assignment
\[
   \Fclass \colon L \;\longmapsto\; \Cl(L),
\]
which we call the \emph{constellation map}. An earlier draft of this
paper used the word ``functor'' for $\Fclass$ informally; this revision
reserves that word for the precise categorical setting given in
\Cref{def:F-functor} below and otherwise speaks of the
\emph{constellation map}. In the body of the paper $\Fclass$ is treated
as a map between sets of objects (scalar local models with one dominant
Newton edge on the source side, finite multisets in $\mathbb{C}$ up to
global rotation on the target). \Cref{def:F-functor} promotes this to
a functor between two small categories whose morphisms are formal gauge
equivalences preserving the designated edge (source) and rotation-
compatible bijections of multisets (target); no statement in the paper
depends on the morphism part of $\Fclass$, but the formulation pins
down what we mean by ``the fibres of $\Fclass$''.

\begin{definition}[The category $\mathfrak{L}$ and the constellation
functor]\label{def:F-functor}
Let $\mathfrak{L}$ be the category whose
\begin{itemize}[leftmargin=2em]
\item \emph{objects} are pairs $(L, \mathfrak{e})$, where $L$ is a
scalar linear operator on $\mathbb{C}[x]$ in $\{\opD, \opS, \opT\}$
with exactly one irregular singularity at $x = \infty$ and
$\mathfrak{e}$ is a designated dominant edge of $\NL$ at infinity;
\item \emph{morphisms} $(L, \mathfrak{e}) \to (L', \mathfrak{e}')$
are formal gauge equivalences $\widehat{M}(L) \xrightarrow{\sim}
\widehat{M}(L')$ over $\mathbb{C}((x^{-1}))$ that send
$\mathfrak{e}$ to $\mathfrak{e}'$ under the induced action on Newton
polygons.
\end{itemize}
Let $\mathfrak{C}$ be the category whose
\begin{itemize}[leftmargin=2em]
\item \emph{objects} are finite multisets $C \subset \mathbb{C}$
taken up to the global rotation equivalence
$C \sim e^{i\theta} C$ for $\theta \in \mathbb{R}/2\pi\mathbb{Z}$;
\item \emph{morphisms} are multiset bijections compatible with the
rotation action.
\end{itemize}
The assignment $(L, \mathfrak{e}) \mapsto \Cl(L)$, extended to morphisms
by sending a formal gauge equivalence to the induced bijection of
roots of $\chi$, defines a functor
\[
   \Fclass \colon \mathfrak{L} \;\longrightarrow\; \mathfrak{C}.
\]
\end{definition}

\begin{remark}\label{rem:F-functoriality}
Functoriality is immediate from the construction of $\chi$ from the
leading-coefficient data along $\mathfrak{e}$: a formal gauge
equivalence preserving $\mathfrak{e}$ acts on the $\LCinf$-data of
$\mathfrak{e}$ by a cocycle that preserves $\chi$ up to an overall
rescaling of $c$, and hence preserves the multiset of roots of $\chi$
up to a global rotation. We will use ``$\Fclass$'' and ``the
constellation map'' interchangeably in the remainder of the paper.
Whenever only the underlying object map matters, the reader may
forget \Cref{def:F-functor} without loss.
\end{remark}

\subsection{The structural law}\label{ssec:intro-law}

The central observation, proved in \Cref{thm:main}, is the algebraic
identity
\[
   \chi(c) \;=\; \chi_w(c^q),
\]
where $\chi$ is the WKB characteristic polynomial built from the
$\LCinf$-data along the dominant edge and $\chi_w \in \mathbb{C}[w]$
is the \emph{inner polynomial} of degree $r$. Solving $\chi = 0$
amounts to solving $\chi_w$ in the variable $w = c^q$ and pulling
back along $c \mapsto c^q$; the constellation decomposes canonically
as a $\mu_q$-equivariant union of $r$ fibres over $\Roots(\chi_w)$.
This factorisation is at once an obstruction (no multiset outside the
image of $\Fclass$ admits it) and a recipe (every multiset in the
image is so built). \Cref{sec:main} states the law as
\Cref{thm:main} with two corollaries (ring decomposition; Vieta
monodromy) and a remark on the degeneracy strata.

\subsection{$\chi_w$ as the edge-indicial polynomial}
\label{ssec:intro-chi-w}

The polynomial $\chi_w$ is not new content. It is exactly the
\emph{edge-indicial polynomial} of the dominant edge $\mathfrak{e}$
in the variable $w = c^q$, the form that appears in the classical
leading-balance analysis of irregular singularities (Wasow
\cite{wasow}, Sibuya \cite{sibuya}, Sabbah
\cite{sabbah-isomon}). Setting $w = c^q$, the equation $\chi_w(w) = 0$
is the standard requirement that the WKB ansatz produce a
non-trivial formal solution along $\mathfrak{e}$. The novelty of
\Cref{thm:main} is therefore not the existence of $\chi_w$ but the
clean identification
\[
   \chi(c) \;=\; \chi_w(c^q),
\]
together with the consequence that the $\mu_q$-equivariant fibre
structure of $\Cl(L)$ becomes explicit, and the slope denominator $q$
appears as a \emph{pullback index} rather than as a feature of the
leading balance to be rediscovered case by case.

\subsection{From a law to a classification}
\label{ssec:intro-classification}

The structural law turns the classification problem into a
combinatorial one. Up to global rotation, a constellation is recovered
from $(p, q, r)$ together with the modulus-with-multiplicity partition
$\pi$ of the roots of $\chi_w$. A computational sweep, recorded in the
companion dataset \cite{wkb_dataset_v1}, enumerates these classes for
slopes $-p/q \in \{1/2, 1/3, 1/4, 2/3, 3/2\}$ and ranks $r \in
\{1, 2, 3\}$ across three equation types: there are exactly $23$
classes in the single-edge case. Once multi-edge polygons, nested
ramification, and the prime-$q$ arithmetic stratification are recorded,
the atlas refines to $39$ classes. These populate the open stratum
and four named degenerate strata (equal modulus; resonant multiplicity;
zero root; accidental symmetry); \Cref{sec:classes} catalogues them.

The classification is extended in several directions:
\begin{itemize}[leftmargin=*]
\item beyond the single edge: per-edge factorisation along the slope
      filtration (\Cref{thm:M1}, \Cref{sec:multiedge});
\item beyond a single $\mu_q$-orbit: nested ramification with iterated
      pullbacks (\Cref{thm:N1}, \Cref{sec:multiedge});
\item beyond squarefree $\chi_w$: degeneracy strata and a prime-$q$
      Kummer dichotomy (\Cref{thm:D2}, \Cref{sec:degeneracy});
\item beyond scalar operators: block-diagonal systems and a
      symmetric-power target (\Cref{thm:S1}, \Cref{sec:systems});
\item beyond the discrete invariant: an analytic refinement attaching
      anti-Stokes directions to the differences $\Dij = c_i - c_j$
      (\Cref{thm:AT1,thm:AT3,thm:AT5,thm:AT7}, \Cref{sec:analytic}).
\end{itemize}

\subsection{Scope of the analytic refinement}
\label{ssec:intro-analytic-scope}

The analytic theorems of \Cref{sec:analytic} attach Stokes-side
information to the constellation. The reader should bear in mind a
clear division of labour:
\begin{itemize}[leftmargin=2em]
\item \emph{Classical content.} The Levelt--Turrittin slope
decomposition of a formal connection on a ramified cover of index $q$;
the existence of formal exponential factors $Q_i(t) = c_i t^p +
(\text{lower})$; the principle that anti-Stokes rays are governed by
the leading term of $Q_i - Q_j$; and the slope-graded structure of
Stokes data for multi-edge connections are classical
(Sibuya \cite{sibuya}, Wasow \cite{wasow}, Sabbah
\cite{sabbah-polarizable}).
\item \emph{Algebraic refinement.} What is new in the present paper
is the explicit linear form for the anti-Stokes condition in
\Cref{thm:AT1},
\[
   \arg \Dij \;+\; p\phi \;\equiv\; \tfrac{\pi}{2} \pmod \pi;
\]
the observation that nested factorisations of $\chi_w$ induce an
additional $\mu_{q'}$-action on the discrete part of the wild data
\emph{without} changing the ramification index $q$ (\Cref{thm:AT3});
the consequent $\mu_q/\mu_{q'}$ symmetry bookkeeping; and the
$p$-visibility statement \Cref{thm:AT7}, by which the slope numerator
$p$---invisible to $\Cl(L)$---is visible to the cover-side ray pattern.
\item \emph{Equivariance, not Stokes matrices.} \Cref{thm:AT3} asserts
$\mu_{q'}$-equivariance of an irregular-type parametrisation; it is
not a statement about wild character variety strata in the moduli-
theoretic sense of Boalch \cite{boalch-stokes,boalch-wildhodge}.
We are explicit about this distinction.
\end{itemize}

\subsection{What is and is not classified}\label{ssec:intro-scope}

The object classified in this paper is the image of $\Fclass$ over
the single-edge ($23$-class) and the extended single-edge / multi-edge
/ nested-ramification / system / arithmetic strata ($39$-class). We
do \emph{not} classify operators. The constellation is a coarsening of
the irregular-type data of Levelt--Turrittin, and the kernel of this
coarsening---operators with identical constellation but distinct
Stokes structures---is the complementary problem solved (in
appropriate categories) by wild character varieties (Boalch
\cite{boalch-stokes,boalch-wildhodge}). \Cref{sec:related} positions
the present work alongside neighbouring programmes: ADE
classifications of singularities; Stokes-graph and spectral-network
classifications; the formal classification of irregular connections
(Levelt--Turrittin, Sabbah \cite{sabbah-isomon,sabbah-polarizable},
Mochizuki \cite{mochizuki-asymptotic}); and wild character varieties.
The constellation classification is orthogonal to each.

\subsection{Roadmap}\label{ssec:intro-organisation}

\Cref{sec:prelim} fixes notation, recalls Newton polygons and the WKB
symbol, and includes a short notation table.
\Cref{sec:main} proves the structural law (\Cref{thm:main}).
\Cref{sec:multiedge} extends to multi-edge polygons via the slope-
filtration per-edge factorisation (\Cref{thm:M1}) and to nested
ramification (\Cref{thm:N1}).
\Cref{sec:degeneracy} treats the four degeneracy strata and proves the
prime-$q$ Kummer dichotomy (\Cref{thm:D2}), separating the algebraic
statement (Kummer theory over $\mathbb{Q}(\zeta_q)$) from the
empirical bounded-search corroboration.
\Cref{sec:systems} treats systems (\Cref{thm:S1}).
\Cref{sec:analytic} states and proves the four analytic-lifting
theorems (\Cref{thm:AT1,thm:AT3,thm:AT5,thm:AT7}), with explicit
hypotheses (OS1)--(OS5) collected in \Cref{ssec:OS}.
\Cref{sec:classes} presents the $23$-class atlas, the $39$-class
refinement, the taxonomy diagram (\Cref{fig:taxonomy}), and the
constellation zoo (\Cref{fig:zoo}).
\Cref{sec:related} positions the work alongside its neighbours,
\Cref{sec:vision} closes with future directions, and
\Cref{sec:reproducibility} contains a reproducibility note.

\bigskip
The reader interested only in the structural law may read
\Cref{sec:main} immediately after \Cref{sec:prelim}; the analytic
theorems of \Cref{sec:analytic} depend logically on
\Cref{sec:main,sec:multiedge} but may be read in any internal order.
